% ID: SEL-D0049-OB-E02
% ESTADO: cotejada. Segunda pasada visual y borrador compilado revisados.
% Fuente: Selectividad/sel_2013_matematicas_aplicadas/junio/reserva b/EXAMEN 1-ANDALUCÍA.pdf
% SHA-256: ec6fd00d0daee547aa6128f3f6e1099278ee551f7f80917b331a1772764317b3
% Página: 2; opción B; ejercicio original 2.

Sea la función $f(x)=\dfrac{1}{3}x^3+\dfrac{1}{2}x^2-2x+3$.
\begin{enumerate}[label=\alph*)]
\item \textbf{(1 punto)} Determine sus máximos y mínimos relativos.
\item \textbf{(1 punto)} Consideremos la función $g(x)=f'(x)$. Calcule la ecuación de la recta tangente a la gráfica de la función $g(x)$, en el punto de abscisa $x=2$.
\item \textbf{(0.5 puntos)} Dibuje la gráfica de $g(x)$ y de la recta tangente calculada en b).
\end{enumerate}
